\section{Introduction}
\label{sec:intro}

Hybrid architectures that interleave selective state-space mixers~\citep{gu2023mamba} with attention~\citep{vaswani2017attention} have emerged as a strong alternative to pure-attention transformers, combining the linear-time long-range modelling of structured state-space models (SSMs) with the precise content-based retrieval that attention provides~\citep{lieber2024jamba,glorioso2024zamba,park2024empirical}. These hybrids retain attention's expressivity at the locations where it matters most, while relegating the bulk of the sequence mixing to the more parameter-efficient SSM path.

We argue that two further architectural modifications are natural and underexplored within this hybrid framing. The first concerns the SSM's \emph{internal} structure. The selective state h_t \in \mathbb{R}^{\dinner \times \dstate} of a Mamba block is best read as $\dinner$ parallel channels, each carrying a $\dstate$-dimensional memory. Within each channel the $\dstate$ slots evolve completely independently --- they only ever interact through the linear readout $y_t = h_t C_t^\top$. Yet the readout vector $C_t$ is already query-like in spirit: it is computed from the input and selects which slots to attend to. We propose to take this selection seriously and replace the linear contraction with full content-dependent attention over the $\dstate$ slots, which we term \emph{intracellular attention}.

The second modification concerns the residual stream \emph{between} blocks. Activation steering and representation engineering have shown that a small number of linear directions in a transformer's residual stream govern semantically meaningful aspects of its behaviour, such as truthfulness, tone, and refusal~\citep{zou2023representation,turner2023activation,rimsky2024steering}. These directions are typically applied at inference time, manually, and post-hoc. We instead extract such directions from a partially-trained base model via document-style contrast pairs and re-inject them \emph{during pretraining} through a learned routing layer that selects a token-conditional convex combination of directions, applies a learned per-direction magnitude, and gates the result through a zero-initialised scalar. We call this construction \emph{hormone routing}: behavioural ``hormones'' that the model learns to release, suppress, and reinterpret as it trains.

Both mechanisms share a common design discipline. Each adds fewer than 1\% to the base model's parameter count. Each is gated such that, at initialisation, the modification is an exact no-op and the surrounding hybrid backbone is recovered bit-for-bit. Each is structurally orthogonal to the other --- intracellular attention modifies how an SSM block reads from its own state, hormone routing modifies what a downstream block sees in the residual --- and we show empirically that they compose without interference. Each can therefore be ablated in isolation.

Our contributions are as follows.
\begin{itemize}
    \item We introduce \textbf{intracellular attention}, a content-dependent slot-mixing mechanism placed inside the SSM scan that lets the $\dstate$ state slots attend to each other before the linear readout. The mechanism preserves the SSM's recurrent form and adds parameters proportional only to a small bottleneck dimension $\dattn \ll \dinner$.
    \item We introduce \textbf{hormone routing}, an architectural component that re-injects extracted behavioural direction vectors into the residual stream during pretraining via a learned per-token router with per-direction magnitudes and a zero-initialised output gate.
    \item We propose a \textbf{document-style contrast-pair extraction protocol} that recovers behavioural directions from the base model itself --- without instruction-tuning data or annotated preferences --- and discuss the methodological reasons for preferring it to instruction-style pairs in the pretraining setting.
    \item We pretrain a 125M-parameter hybrid Mamba/MQA backbone on FineWeb-Edu and present a controlled ablation study isolating each contribution, reporting validation perplexity and zero-shot accuracy on HellaSwag, LAMBADA, ARC-Easy, ARC-Challenge, and PIQA across seven variants spanning both contributions and their composition.
\end{itemize}

The rest of the paper is structured as follows. Section~\ref{sec:related} situates the work against the SSM, hybrid, and representation-engineering literatures. Section~\ref{sec:arch} gives the hybrid backbone. Sections~\ref{sec:intracellular} and~\ref{sec:hormone} present the two contributions in turn. Sections~\ref{sec:experiments} and~\ref{sec:analysis} report the empirical results and ablations.
